% --- Archon physics formalization source begin ---
% archon:physics
% archon:covers PhyXMiniProblems/problem_phyx_mini_0218.lean
% archon:source-report reports/phyx_mini/problem_phyx_mini_0218.source.json

\chapter{Physics problem phyx\_mini\_0218}
\label{ch:PhyXMiniProblems_problem_phyx_mini_0218}

\paragraph{Problem source.}
\#\# Physical scenario

When a player's finger presses a guitar string down onto a fret, the length of the vibrating portion of the string is shortened, thereby increasing the string's fundamental frequency (see figure). The string's tension and mass per unit length remain unchanged.The unfingered length of the string is \textbackslash{}( \textbackslash{}ell = 75.0\textbackslash{},\textbackslash{}mathrm\{cm\} \textbackslash{}). If each fret raises the pitch of the fundamental by one musical note compared to the neighboring fret. On the equally tempered chromatic scale, the ratio of frequencies of neighboring notes is \textbackslash{}( 2\textasciicircum{}\{1/12\} \textbackslash{}).

\#\# Question

Determine the positions \textbackslash{}( x \textbackslash{}) of the first fret.

\#\# Answer choices

- A: A: 8.2cm

- B: B: 6.2cm

- C: C: 5.2cm

- D: D: 4.2cm

\#\# Figure caption (auxiliary; use the image as primary evidence)

The image depicts a guitar with a hand positioned over the strings. The guitar is oriented horizontally, showing the body, neck, strings, and tuning pegs. A hand is illustrated in a position to pluck or interact with the strings near the neck. 

There are two measurements indicated:
1. The total length of the guitar is annotated as \textbackslash{}( l = 75.0 \textbackslash{}, \textbackslash{}text\{cm\} \textbackslash{}).
2. Another distance along the neck is marked as \textbackslash{}( x \textbackslash{}).

The measurements have arrows pointing to specific parts of the guitar, illustrating the areas being measured. The strings are visible and run the length of the neck, over the soundhole and bridge.

\#\# Dataset metadata

Resonance and Harmonics; Physical Model Grounding Reasoning, Numerical Reasoning

\paragraph{Recorded answer/context.}
D: D: 4.2cm

\paragraph{Figure/image path.}
/root/proposal\_for\_physic/science-mango/phyx\_mini\_run/phyx\_data/test\_image/218.png

\paragraph{Formalization target.}
create a compiling Lean file with sorry bodies at `PhyXMiniProblems/problem\_phyx\_mini\_0218.lean`. The Lean declarations must preserve the physical quantities, dimensions or dimensional roles, figure labels, governing-law hypotheses, and final relation expressed by this problem.
Use Mathlib/Physlib names found through LeanExplore where available. If a physics API is missing, introduce faithful local abstractions rather than scalar placeholder aliases.

\begin{theorem}[Physics formalization target]
\label{thm:physics:phyx_mini_0218:target}
\lean{PhyXMiniProblems.ProblemPhyXMini0218.firstFretPosition_matches_recordedAnswerD}
\uses{def:physics:phyx-mini-0218:phyxminiproblems-problemphyxmini0218-lengthincentimeters, def:physics:phyx-mini-0218:phyxminiproblems-problemphyxmini0218-guitarstringsetup, def:physics:phyx-mini-0218:phyxminiproblems-problemphyxmini0218-matchesproblemandprimaryfigure, def:physics:phyx-mini-0218:phyxminiproblems-problemphyxmini0218-hasphysicalstringparameters, def:physics:phyx-mini-0218:phyxminiproblems-problemphyxmini0218-satisfiesstretchedstringfundamentallaw, def:physics:phyx-mini-0218:phyxminiproblems-problemphyxmini0218-equaltemperedsemitoneratio, def:physics:phyx-mini-0218:phyxminiproblems-problemphyxmini0218-raisespitchbyoneequaltemperedsemitone, def:physics:phyx-mini-0218:phyxminiproblems-problemphyxmini0218-recordedanswerchoice, def:physics:phyx-mini-0218:phyxminiproblems-problemphyxmini0218-matchesanswerchoice}
The assigned autoformalize agent should translate this physics problem into Lean declarations in the covered file, with theorem and lemma proof bodies written as `by sorry`.
\end{theorem}
\begin{proof}
This is an autoformalization task, not a proof task. Produce statements that can later be proved by the physics prover without weakening the physical model.
\end{proof}

% --- Archon declaration topology begin ---
\section{Declaration topology}
\begin{definition}[Length Quantity]
  \label{def:physics:phyx-mini-0218:phyxminiproblems-problemphyxmini0218-lengthquantity}
  \lean{PhyXMiniProblems.ProblemPhyXMini0218.LengthQuantity}
  A nonnegative physical length, used for ℓ, x, and vibrating lengths
\end{definition}
\begin{proof}
  This is the declared model definition of Length Quantity.
\end{proof}
\begin{definition}[Frequency Quantity]
  \label{def:physics:phyx-mini-0218:phyxminiproblems-problemphyxmini0218-frequencyquantity}
  \lean{PhyXMiniProblems.ProblemPhyXMini0218.FrequencyQuantity}
  A nonnegative ordinary frequency, with inverse-time dimension
\end{definition}
\begin{proof}
  This is the declared model definition of Frequency Quantity.
\end{proof}
\begin{definition}[Tension Quantity]
  \label{def:physics:phyx-mini-0218:phyxminiproblems-problemphyxmini0218-tensionquantity}
  \lean{PhyXMiniProblems.ProblemPhyXMini0218.TensionQuantity}
  A nonnegative string tension, carrying the physical dimension of force
\end{definition}
\begin{proof}
  This is the declared model definition of Tension Quantity.
\end{proof}
\begin{definition}[Linear Mass Density Quantity]
  \label{def:physics:phyx-mini-0218:phyxminiproblems-problemphyxmini0218-linearmassdensityquantity}
  \lean{PhyXMiniProblems.ProblemPhyXMini0218.LinearMassDensityQuantity}
  A nonnegative mass per unit length of guitar string
\end{definition}
\begin{proof}
  This is the declared model definition of Linear Mass Density Quantity.
\end{proof}
\begin{definition}[length In Meters]
  \label{def:physics:phyx-mini-0218:phyxminiproblems-problemphyxmini0218-lengthinmeters}
  \lean{PhyXMiniProblems.ProblemPhyXMini0218.lengthInMeters}
  \uses{def:physics:phyx-mini-0218:phyxminiproblems-problemphyxmini0218-lengthquantity}
  Metre readout of a physical length in coherent SI units
\end{definition}
\begin{proof}
  This is the declared model definition of length In Meters.
\end{proof}
\begin{definition}[length In Centimeters]
  \label{def:physics:phyx-mini-0218:phyxminiproblems-problemphyxmini0218-lengthincentimeters}
  \lean{PhyXMiniProblems.ProblemPhyXMini0218.lengthInCentimeters}
  \uses{def:physics:phyx-mini-0218:phyxminiproblems-problemphyxmini0218-lengthquantity, def:physics:phyx-mini-0218:phyxminiproblems-problemphyxmini0218-lengthinmeters}
  Centimetre readout derived from the coherent SI metre readout
\end{definition}
\begin{proof}
  This is the declared model definition of length In Centimeters.
\end{proof}
\begin{definition}[frequency In Hertz]
  \label{def:physics:phyx-mini-0218:phyxminiproblems-problemphyxmini0218-frequencyinhertz}
  \lean{PhyXMiniProblems.ProblemPhyXMini0218.frequencyInHertz}
  \uses{def:physics:phyx-mini-0218:phyxminiproblems-problemphyxmini0218-frequencyquantity}
  Hertz readout of a physical frequency
\end{definition}
\begin{proof}
  This is the declared model definition of frequency In Hertz.
\end{proof}
\begin{definition}[tension In Newtons]
  \label{def:physics:phyx-mini-0218:phyxminiproblems-problemphyxmini0218-tensioninnewtons}
  \lean{PhyXMiniProblems.ProblemPhyXMini0218.tensionInNewtons}
  \uses{def:physics:phyx-mini-0218:phyxminiproblems-problemphyxmini0218-tensionquantity}
  Newton readout of a physical string tension
\end{definition}
\begin{proof}
  This is the declared model definition of tension In Newtons.
\end{proof}
\begin{definition}[linear Mass Density In Kilograms Per Meter]
  \label{def:physics:phyx-mini-0218:phyxminiproblems-problemphyxmini0218-linearmassdensityinkilogramspermeter}
  \lean{PhyXMiniProblems.ProblemPhyXMini0218.linearMassDensityInKilogramsPerMeter}
  \uses{def:physics:phyx-mini-0218:phyxminiproblems-problemphyxmini0218-linearmassdensityquantity}
  Kilogram-per-metre readout of a physical linear mass density
\end{definition}
\begin{proof}
  This is the declared model definition of linear Mass Density In Kilograms Per Meter.
\end{proof}
\begin{definition}[String Configuration]
  \label{def:physics:phyx-mini-0218:phyxminiproblems-problemphyxmini0218-stringconfiguration}
  \lean{PhyXMiniProblems.ProblemPhyXMini0218.StringConfiguration}
  The two string configurations compared in the problem
\end{definition}
\begin{proof}
  This is the declared model definition of String Configuration.
\end{proof}
\begin{definition}[String Landmark]
  \label{def:physics:phyx-mini-0218:phyxminiproblems-problemphyxmini0218-stringlandmark}
  \lean{PhyXMiniProblems.ProblemPhyXMini0218.StringLandmark}
  Ordered landmarks along the string in the primary figure
\end{definition}
\begin{proof}
  This is the declared model definition of String Landmark.
\end{proof}
\begin{definition}[Vibrating Segment End]
  \label{def:physics:phyx-mini-0218:phyxminiproblems-problemphyxmini0218-vibratingsegmentend}
  \lean{PhyXMiniProblems.ProblemPhyXMini0218.VibratingSegmentEnd}
  The two endpoints of the portion of string that is free to vibrate
\end{definition}
\begin{proof}
  This is the declared model definition of Vibrating Segment End.
\end{proof}
\begin{definition}[Guitar Orientation]
  \label{def:physics:phyx-mini-0218:phyxminiproblems-problemphyxmini0218-guitarorientation}
  \lean{PhyXMiniProblems.ProblemPhyXMini0218.GuitarOrientation}
  The horizontal orientation in which the guitar is drawn
\end{definition}
\begin{proof}
  This is the declared model definition of Guitar Orientation.
\end{proof}
\begin{definition}[Guitar String Setup]
  \label{def:physics:phyx-mini-0218:phyxminiproblems-problemphyxmini0218-guitarstringsetup}
  \lean{PhyXMiniProblems.ProblemPhyXMini0218.GuitarStringSetup}
  \uses{def:physics:phyx-mini-0218:phyxminiproblems-problemphyxmini0218-lengthquantity, def:physics:phyx-mini-0218:phyxminiproblems-problemphyxmini0218-frequencyquantity, def:physics:phyx-mini-0218:phyxminiproblems-problemphyxmini0218-tensionquantity, def:physics:phyx-mini-0218:phyxminiproblems-problemphyxmini0218-linearmassdensityquantity, def:physics:phyx-mini-0218:phyxminiproblems-problemphyxmini0218-stringconfiguration, def:physics:phyx-mini-0218:phyxminiproblems-problemphyxmini0218-stringlandmark, def:physics:phyx-mini-0218:phyxminiproblems-problemphyxmini0218-vibratingsegmentend, def:physics:phyx-mini-0218:phyxminiproblems-problemphyxmini0218-guitarorientation}
  Independent physical quantities and geometric labels for the guitar string. firstFretPositionFromNut is the quantity labeled x in the figure. It is left unconstrained here: in particular, neither 4.2 cm nor the exact equal-temperament expression is built into this setup
\end{definition}
\begin{proof}
  This is the declared model definition of Guitar String Setup.
\end{proof}
\begin{definition}[Matches Problem And Primary Figure]
  \label{def:physics:phyx-mini-0218:phyxminiproblems-problemphyxmini0218-matchesproblemandprimaryfigure}
  \lean{PhyXMiniProblems.ProblemPhyXMini0218.MatchesProblemAndPrimaryFigure}
  \uses{def:physics:phyx-mini-0218:phyxminiproblems-problemphyxmini0218-lengthincentimeters, def:physics:phyx-mini-0218:phyxminiproblems-problemphyxmini0218-guitarstringsetup}
  The numerical and geometric information stated in the problem and read from the primary figure. The first-fret distance x plus the remaining vibrating length equals the original bridge-to-nut length ℓ. The unchanged-tension and unchanged-density fields are stated data, not the requested fret-position conclusion
\end{definition}
\begin{proof}
  This is the declared model definition of Matches Problem And Primary Figure.
\end{proof}
\begin{definition}[Has Physical String Parameters]
  \label{def:physics:phyx-mini-0218:phyxminiproblems-problemphyxmini0218-hasphysicalstringparameters}
  \lean{PhyXMiniProblems.ProblemPhyXMini0218.HasPhysicalStringParameters}
  \uses{def:physics:phyx-mini-0218:phyxminiproblems-problemphyxmini0218-lengthinmeters, def:physics:phyx-mini-0218:phyxminiproblems-problemphyxmini0218-frequencyinhertz, def:physics:phyx-mini-0218:phyxminiproblems-problemphyxmini0218-tensioninnewtons, def:physics:phyx-mini-0218:phyxminiproblems-problemphyxmini0218-linearmassdensityinkilogramspermeter, def:physics:phyx-mini-0218:phyxminiproblems-problemphyxmini0218-guitarstringsetup}
  Positivity and nondegeneracy conditions for an ordinary taut string
\end{definition}
\begin{proof}
  This is the declared model definition of Has Physical String Parameters.
\end{proof}
\begin{definition}[Satisfies Stretched String Fundamental Law]
  \label{def:physics:phyx-mini-0218:phyxminiproblems-problemphyxmini0218-satisfiesstretchedstringfundamentallaw}
  \lean{PhyXMiniProblems.ProblemPhyXMini0218.SatisfiesStretchedStringFundamentalLaw}
  \uses{def:physics:phyx-mini-0218:phyxminiproblems-problemphyxmini0218-lengthinmeters, def:physics:phyx-mini-0218:phyxminiproblems-problemphyxmini0218-frequencyinhertz, def:physics:phyx-mini-0218:phyxminiproblems-problemphyxmini0218-tensioninnewtons, def:physics:phyx-mini-0218:phyxminiproblems-problemphyxmini0218-linearmassdensityinkilogramspermeter, def:physics:phyx-mini-0218:phyxminiproblems-problemphyxmini0218-guitarstringsetup}
  The fixed-end fundamental-mode law for each configuration, f₁ = (1 / (2 L)) sqrt(T / μ). This is a governing law relating the independent physical quantities. It does not mention the value or formula requested for the fret position
\end{definition}
\begin{proof}
  This is the declared model definition of Satisfies Stretched String Fundamental Law.
\end{proof}
\begin{definition}[equal Tempered Semitone Ratio]
  \label{def:physics:phyx-mini-0218:phyxminiproblems-problemphyxmini0218-equaltemperedsemitoneratio}
  \lean{PhyXMiniProblems.ProblemPhyXMini0218.equalTemperedSemitoneRatio}
  The dimensionless ratio between neighboring equal-tempered notes
\end{definition}
\begin{proof}
  This is the declared model definition of equal Tempered Semitone Ratio.
\end{proof}
\begin{definition}[Raises Pitch By One Equal Tempered Semitone]
  \label{def:physics:phyx-mini-0218:phyxminiproblems-problemphyxmini0218-raisespitchbyoneequaltemperedsemitone}
  \lean{PhyXMiniProblems.ProblemPhyXMini0218.RaisesPitchByOneEqualTemperedSemitone}
  \uses{def:physics:phyx-mini-0218:phyxminiproblems-problemphyxmini0218-frequencyinhertz, def:physics:phyx-mini-0218:phyxminiproblems-problemphyxmini0218-guitarstringsetup, def:physics:phyx-mini-0218:phyxminiproblems-problemphyxmini0218-equaltemperedsemitoneratio}
  The stated musical relation: fingering the first fret raises the fundamental by one equal-tempered semitone. This gives a frequency ratio, not the unknown geometric position x
\end{definition}
\begin{proof}
  This is the declared model definition of Raises Pitch By One Equal Tempered Semitone.
\end{proof}
\begin{lemma}[fingered Vibrating Length ratio]
  \label{lem:physics:phyx-mini-0218:phyxminiproblems-problemphyxmini0218-fingeredvibratinglength-ratio}
  \lean{PhyXMiniProblems.ProblemPhyXMini0218.fingeredVibratingLength_ratio}
  \uses{def:physics:phyx-mini-0218:phyxminiproblems-problemphyxmini0218-lengthinmeters, def:physics:phyx-mini-0218:phyxminiproblems-problemphyxmini0218-guitarstringsetup, def:physics:phyx-mini-0218:phyxminiproblems-problemphyxmini0218-matchesproblemandprimaryfigure, def:physics:phyx-mini-0218:phyxminiproblems-problemphyxmini0218-hasphysicalstringparameters, def:physics:phyx-mini-0218:phyxminiproblems-problemphyxmini0218-satisfiesstretchedstringfundamentallaw, def:physics:phyx-mini-0218:phyxminiproblems-problemphyxmini0218-equaltemperedsemitoneratio, def:physics:phyx-mini-0218:phyxminiproblems-problemphyxmini0218-raisespitchbyoneequaltemperedsemitone}
  For fixed tension and linear density, the stretched-string law and the given semitone ratio imply that the fingered vibrating length is the unfingered length divided by 2\textasciicircum{}(1/12). This lemma remains a derived conclusion
\end{lemma}
\begin{proof}
  The conclusion follows from the governing relations and figure readouts named in the statement of fingered Vibrating Length ratio.
\end{proof}
\begin{lemma}[first Fret Position exact Formula]
  \label{lem:physics:phyx-mini-0218:phyxminiproblems-problemphyxmini0218-firstfretposition-exactformula}
  \lean{PhyXMiniProblems.ProblemPhyXMini0218.firstFretPosition_exactFormula}
  \uses{def:physics:phyx-mini-0218:phyxminiproblems-problemphyxmini0218-lengthincentimeters, def:physics:phyx-mini-0218:phyxminiproblems-problemphyxmini0218-guitarstringsetup, def:physics:phyx-mini-0218:phyxminiproblems-problemphyxmini0218-matchesproblemandprimaryfigure, def:physics:phyx-mini-0218:phyxminiproblems-problemphyxmini0218-hasphysicalstringparameters, def:physics:phyx-mini-0218:phyxminiproblems-problemphyxmini0218-satisfiesstretchedstringfundamentallaw, def:physics:phyx-mini-0218:phyxminiproblems-problemphyxmini0218-equaltemperedsemitoneratio, def:physics:phyx-mini-0218:phyxminiproblems-problemphyxmini0218-raisespitchbyoneequaltemperedsemitone}
  The corresponding exact formula for the figure's nut-to-first-fret distance. It is stated before rounding to any displayed answer
\end{lemma}
\begin{proof}
  The conclusion follows from the governing relations and figure readouts named in the statement of first Fret Position exact Formula.
\end{proof}
\begin{definition}[Answer Choice]
  \label{def:physics:phyx-mini-0218:phyxminiproblems-problemphyxmini0218-answerchoice}
  \lean{PhyXMiniProblems.ProblemPhyXMini0218.AnswerChoice}
  Labels of the four answer choices printed with the problem
\end{definition}
\begin{proof}
  This is the declared model definition of Answer Choice.
\end{proof}
\begin{definition}[centimeters]
  \label{def:physics:phyx-mini-0218:phyxminiproblems-problemphyxmini0218-answerchoice-centimeters}
  \lean{PhyXMiniProblems.ProblemPhyXMini0218.AnswerChoice.centimeters}
  \uses{def:physics:phyx-mini-0218:phyxminiproblems-problemphyxmini0218-answerchoice}
  Centimetre value printed beside each answer label
\end{definition}
\begin{proof}
  This is the declared model definition of centimeters.
\end{proof}
\begin{definition}[recorded Answer Choice]
  \label{def:physics:phyx-mini-0218:phyxminiproblems-problemphyxmini0218-recordedanswerchoice}
  \lean{PhyXMiniProblems.ProblemPhyXMini0218.recordedAnswerChoice}
  \uses{def:physics:phyx-mini-0218:phyxminiproblems-problemphyxmini0218-answerchoice}
  The answer label recorded in the supplied dataset metadata
\end{definition}
\begin{proof}
  This is the declared model definition of recorded Answer Choice.
\end{proof}
\begin{definition}[Matches Answer Choice]
  \label{def:physics:phyx-mini-0218:phyxminiproblems-problemphyxmini0218-matchesanswerchoice}
  \lean{PhyXMiniProblems.ProblemPhyXMini0218.MatchesAnswerChoice}
  \uses{def:physics:phyx-mini-0218:phyxminiproblems-problemphyxmini0218-lengthquantity, def:physics:phyx-mini-0218:phyxminiproblems-problemphyxmini0218-lengthincentimeters, def:physics:phyx-mini-0218:phyxminiproblems-problemphyxmini0218-answerchoice}
  A physical fret position agrees with a displayed tenth-centimetre answer when its centimetre readout lies within half of 0.1 cm of that displayed value
\end{definition}
\begin{proof}
  This is the declared model definition of Matches Answer Choice.
\end{proof}
% --- Archon declaration topology end ---
% --- Archon physics formalization source end ---
